\section*{Hoofdstuk \ref{ch:g11:lithosphere-asthenosphere} --- Lithosfeer en asthenosfeer}

\begin{solution}{exo:g11:lithosphere-asthenosphere:1}
De \omterm{def:g11:lithosphere-asthenosphere:lithosphere}{lithosfeer} is de starre buitenste schil van de Aarde, opgebouwd uit de korst
en de bovenste mantel; de \omterm{def:g11:lithosphere-asthenosphere:lithosphere}{asthenosfeer} is de hetere, zwakkere mantel eronder,
die over lange tijdspannen langzaam vloeit. De oceanische korst, de continentale
korst en de bovenste mantel horen bij de \omterm{def:g11:lithosphere-asthenosphere:lithosphere}{lithosfeer}; de buitenkern
niet.
\end{solution}

\begin{solution}{exo:g11:lithosphere-asthenosphere:2}
De ijsberg zakt tot hij zijn eigen gewicht aan water opzij heeft geduwd: het deel
onder water is $910/1025 \approx 0.888$, ongeveer \qty{89}{\percent}.
Het deel erboven is ongeveer \qty{11.2}{\percent}, dus de totale hoogte is
$20/0.112 \approx \qty{180}{m}$, waarvan ongeveer \qty{160}{m} onder
water.
\end{solution}

\begin{solution}{exo:g11:lithosphere-asthenosphere:3}
$15 + 26 \times 3 = \qty{93}{\celsius}$. Op \qty{100}{km} zou dezelfde gradiënt
meer dan \qty{2600}{\celsius} geven, ver boven de smelttemperatuur van
mantelgesteente; maar \omterm{def:g11:seismic-structure:body-waves}{S-golven} gaan door de bovenmantel, die dus vast is.
De gradiënt moet met de diepte veel kleiner worden.
\end{solution}

\begin{solution}{exo:g11:lithosphere-asthenosphere:4}
De band van 0 tot \qty{500}{m} boven zeeniveau, met ongeveer
\qty{16}{\percent} van het oppervlak, en de band van 4000 tot \qty{4500}{m}
onder zeeniveau, met ongeveer \qty{11}{\percent} (de band van 4500 tot
\qty{5000}{m} eronder is bijna even groot).
\end{solution}

\begin{solution}{exo:g11:lithosphere-asthenosphere:5}
Bij $P = 0$: \qty{1121}{\celsius}. Bij \qty{3}{GPa}: $-5.104 \times 9 +
132.899 \times 3 + 1120.661 \approx -45.9 + 398.7 + 1120.7 \approx
\qty{1473}{\celsius}$.
\end{solution}

\begin{solution}{exo:g11:lithosphere-asthenosphere:6}
Ongeveer \qty{730}{\celsius} op \qty{50}{km} (\omterm{def:g11:lithosphere-asthenosphere:lithosphere}{lithosfeer}), ongeveer
\qty{1370}{\celsius} op \qty{150}{km} en \qty{1410}{\celsius} op
\qty{250}{km} (\omterm{def:g11:lithosphere-asthenosphere:lithosphere}{asthenosfeer}). De temperatuur stijgt het meest tussen 50 en
\qty{150}{km}, met meer dan \qty{600}{\celsius}, omdat het grootste deel van dat
interval in de \omterm{def:g11:lithosphere-asthenosphere:lithosphere}{lithosfeer} ligt; tussen 150 en \qty{250}{km} stijgt ze maar
met ongeveer \qty{45}{\celsius}.
\end{solution}

\begin{solution}{exo:g11:lithosphere-asthenosphere:7}
Wortel: $3 \times 2830/550 \approx \qty{15.4}{km}$. Extra dikte:
$3 \times 3380/550 \approx \qty{18.4}{km}$, dat wil zeggen de \qty{3}{km} van de
keten plus haar wortel.
\end{solution}

\begin{solution}{exo:g11:lithosphere-asthenosphere:8}
De korst is dikker met $55 - 41.1 = \qty{13.9}{km}$, en $\Delta H = h
\times 3380/550$, dus $h = 13.9 \times 550/3380 \approx \qty{2.3}{km}$.
\end{solution}

\begin{solution}{exo:g11:lithosphere-asthenosphere:9}
Afname: $(4.470 - 4.419)/4.470 \approx \qty{1.1}{\percent}$. Sprong:
$(4.644 - 4.419)/4.419 \approx \qty{5.1}{\percent}$. Met de diepte wordt het
gesteente steeds meer samengedrukt, wat seismische golven sneller maakt; een
afname betekent dat iets anders wint: het gesteente is heet, dicht bij het
smelten, en zwakker.
\end{solution}

\begin{solution}{exo:g11:lithosphere-asthenosphere:10}
(1) $(1345 - 610)/60 \approx \qty{12}{\celsius}$ per kilometer. (2) Dat is
ongeveer de helft van de gradiënt aan het oppervlak. Een deel van de warmte
wordt in de bovenkorst zelf gemaakt; ook die warmte moet door de bovenste lagen
om te ontsnappen, wat de gradiënt dicht bij het oppervlak steiler maakt. Onder de
korst gaat alleen de warmte van dieper door het gesteente, en is de gradiënt kleiner.
\end{solution}

\begin{solution}{exo:g11:lithosphere-asthenosphere:11}
Op \qty{200}{km}: $1300 + 90 = \qty{1390}{\celsius}$; op \qty{400}{km}:
$1300 + 180 = \qty{1480}{\celsius}$. De stijging is \qty{90}{\celsius} over
\qty{200}{km}. Door de \omterm{def:g11:lithosphere-asthenosphere:lithosphere}{lithosfeer} stijgt de temperatuur met
$1345 - 15 = \qty{1330}{\celsius}$ over maar \qty{100}{km}: ongeveer dertig
keer meer per kilometer.
\end{solution}

\begin{solution}{exo:g11:lithosphere-asthenosphere:12}
(1) $118 + 12.3\,z = 1300 + 0.45\,z$ geeft $z = 1182/11.85 \approx
\qty{100}{km}$. (2) $118 + 21.5\,z = 1300 + 0.45\,z$ geeft $z = 1182/21.05
\approx \qty{56}{km}$. (3) Het hetere westen heeft een veel dunnere \omterm{def:g11:lithosphere-asthenosphere:lithosphere}{lithosfeer},
zoals gemeten (ongeveer \qty{40}{km} op het dunste punt); de heel dikke
\omterm{def:g11:lithosphere-asthenosphere:lithosphere}{lithosfeer} van het midden van het continent vraagt een nog kleinere gradiënt.
\end{solution}

\begin{solution}{exo:g11:lithosphere-asthenosphere:13}
(1) Volgens de propositie is $\Delta H = h\,\rho_m/(\rho_m - \rho_c)$, dus $h =
\Delta H\,(\rho_m - \rho_c)/\rho_m$. (2) $550/3380 \approx 0.16$: ongeveer een
zesde van elke extra korst verschijnt als hoogte, vijf zesde als wortel. (3) Voor
de ijsberg is het deel boven het water $(1025 - 910)/1025 \approx 0.11$:
dezelfde regel, (dichtheid van de vloeistof $-$ dichtheid van het drijvende
lichaam)/dichtheid van de vloeistof.
\end{solution}

\begin{solution}{exo:g11:lithosphere-asthenosphere:14}
(1) $1489 - 1345 = \qty{144}{\celsius}$. (2) Een beetje smelt verzamelt zich in
dunne filmpjes tussen de korrels van het gesteente. De korrels zitten dan niet
meer stevig in elkaar vast: het gesteente biedt minder weerstand tegen
afschuiving, dus \omterm{def:g11:seismic-structure:body-waves}{S-golven} vertragen, en het kan veel makkelijker vervormen en vloeien.
\end{solution}

\begin{solution}{exo:g11:lithosphere-asthenosphere:15}
(1) Nee. Het land afvlakken zou de continentale korst iets dunner maken, maar ze
zou nog altijd dik en licht zijn, en nog altijd veel hoger drijven dan de dunne,
dichte korst van de oceaanbodem; de oceanen zouden nog altijd de bekkens tussen
de continenten vullen. (2) Er zijn twee soorten korst, met een verschillende
dikte en dichtheid: een dikke, lichte continentale korst en een dunne, dichte
oceanische korst.
\end{solution}

\begin{solution}{pb:g11:lithosphere-asthenosphere:1}
\textbf{1.} $5.01 - 0.67 = \qty{4.34}{km}$.

\textbf{2.} $(3380 - 2830)/3380 \approx \qty{16}{\percent}$.

\textbf{3.} Het deel $2830/3380 \approx 0.84$ zou eronder liggen; de
ijsberg zakt iets dieper, $910/1025 \approx 0.89$, omdat ijs in dichtheid dichter
bij water ligt dan korst bij mantel.

\textbf{4.} Over de seconden van een seismische golf, of de duur van een
aardbeving, is de mantel een vaste stof; alleen over heel lange tijdspannen
vloeit de \omterm{def:g11:lithosphere-asthenosphere:lithosphere}{asthenosfeer} en laat ze de korst drijven.

\textbf{5.} Onder de diepste \omterm{def:g11:seismic-structure:moho}{Moho} bestaan beide kolommen uit hetzelfde
mantelgesteente op dezelfde diepten, en wegen ze daar evenveel; alle verschillen
liggen boven die diepte en moeten worden meegenomen.

\textbf{6.} Het gemiddelde continent: korst van \qty{0.67}{km} boven tot
\qty{40.4}{km} onder zeeniveau, dan mantel tot \qty{70}{km}. Het
plateau: korst van \qty{5.01}{km} boven zeeniveau tot een \omterm{def:g11:seismic-structure:moho}{Moho} die dieper ligt
met de wortel, dan mantel tot \qty{70}{km}.

\textbf{7.} De kolom van het plateau heeft $h + r$ meer korst, de kolom van het
continent heeft in plaats daarvan mantel over de dikte $r$: $\rho_c\,(h + r) = \rho_m\,r$,
dus $r = h\,\rho_c/(\rho_m - \rho_c)$ en $\Delta H = h + r =
h\,\rho_m/(\rho_m - \rho_c)$.

\textbf{8.} $3380/550 \approx 6.15$.

\textbf{9.} $\Delta H = 4.34 \times 6.15 \approx \qty{26.7}{km}$; de wortel is
$26.7 - 4.34 \approx \qty{22.3}{km}$.

\textbf{10.} De \omterm{def:g11:seismic-structure:moho}{Moho} van het gemiddelde continent ligt $41.1 - 0.67 \approx
\qty{40.4}{km}$ onder zeeniveau; die van het plateau $40.4 + 22.3 \approx
\qty{62.8}{km}$ onder zeeniveau.

\textbf{11.} $41.1 + 26.7 \approx \qty{67.8}{km}$.

\textbf{12.} Een kilometer gesteente boven een vierkante meter heeft een massa van
$\rho \times \qty{1000}{m^3}$: \num{2830} ton korst of \num{3380}
ton mantel. Continent: $41.1 \times 2830 + (70 - 40.43) \times 3380
\approx \num{116300} + \num{99900} \approx \num{216300}$ ton. Plateau:
$67.77 \times 2830 + (70 - 62.76) \times 3380 \approx \num{191800} +
\num{24500} \approx \num{216300}$ ton. De twee zijn gelijk.

\textbf{13.} $41.1 - 0.67 \times 6.15 \approx \qty{37.0}{km}$.

\textbf{14.} $67.8/41.1 \approx 1.65$ keer.

\textbf{15.} $3300/470 \approx 7.02$, $\Delta H \approx 4.34 \times 7.02
\approx \qty{30.5}{km}$, en een korst van ongeveer \qty{71.6}{km}. Een verandering
van ongeveer \qty{2}{\percent} in de dichtheid van de mantel verandert de dikte met
ongeveer \qty{4}{km}: het resultaat hangt merkbaar af van de dichtheden, maar blijft
dicht bij \qty{70}{km}.

\textbf{16.} Het model geeft ongeveer \qty{68}{km}, de seismische golven ongeveer
\qty{70}{km}: ze komen tot op een paar procent overeen.

\textbf{17.} $h = 3.23 - 0.67 = \qty{2.56}{km}$, $\Delta H = 2.56 \times
6.15 \approx \qty{15.7}{km}$, en een korst van ongeveer \qty{56.8}{km}, zo'n
\qty{7}{km} meer dan gemeten.

\textbf{18.} Dunner. Een kolom waarvan het bovenste deel lichter is dan gemiddeld,
weegt minder bij dezelfde dikte, en drijft dus hoger: een lichtere korst heeft
een kleinere wortel nodig om op een gegeven hoogte te staan. De echte korst kan
dus dunner zijn dan die van het model, zoals gemeten.

\textbf{19.} Over lange tijdspannen vloeit de \omterm{def:g11:lithosphere-asthenosphere:lithosphere}{asthenosfeer}: ze wijkt uit onder
een zware kolom en komt terug onder een lichte, tot de kolommen in evenwicht zijn.

\textbf{20.} \textbf{De korst onder Tibet is ongeveer \qty{70}{km} dik}, veel
meer dan de \qty{41}{km} van een gemiddeld continent: het plateau ligt vijf
kilometer hoog omdat zijn dikke, lichte korst op de dichtere mantel drijft,
met een wortel van meer dan \qty{20}{km}.
\end{solution}
